% Generated draft: requires mathematical and source review.
\documentclass[11pt]{article}
\usepackage[margin=25mm]{geometry}
\usepackage{fontspec}
\setmainfont{Latin Modern Roman}
\usepackage{amsmath,amssymb,mathtools,unicode-math}
\setmathfont{Latin Modern Math}
\usepackage{enumitem,needspace,xurl,hyperref}
\setlength{\parindent}{0pt}
\setlength{\parskip}{5pt}
\setlength{\emergencystretch}{3em}
\newcommand{\N}{\mathbb{N}}
\newcommand{\Z}{\mathbb{Z}}
\newcommand{\Q}{\mathbb{Q}}
\newcommand{\R}{\mathbb{R}}
\newcommand{\C}{\mathbb{C}}
\newcommand{\F}{\mathbb{F}}
\newcommand{\A}{\mathbb{A}}
\newcommand{\PP}{\mathbb{P}}
\newcommand{\E}{\mathbb{E}}
\newcommand{\eps}{\varepsilon}
\newcommand{\dd}{\,\mathrm{d}}
\begin{document}
\section*{UnsolvedMath ID: 11000120}
\textbf{Source problem code:} catalogue entry
\subsection{Definitions and mathematical statement}
Let \(\operatorname{Mod}(N)\) be the mapping class group of the surface \(N\) in the source's conventions. The extract does not define \(N\); the research summary discusses nonorientable surfaces, but supplies no genus, punctures, boundary data or mapping-class conventions.

For a group \(G\), \(\operatorname{Aut}(G)\) is its group of abstract automorphisms, \(\operatorname{Inn}(G)\) consists of conjugations \(x\mapsto gxg^{-1}\), and
\[
 \operatorname{Out}(G)=\operatorname{Aut}(G)/\operatorname{Inn}(G),
 \qquad\operatorname{Inn}(G)\cong G/Z(G).
\]
Problem 3.2 on PDF page 98 asks to compute the ``(outer) automorphism group'' of \(\operatorname{Mod}(N)\). Thus the relevant targets are \(\operatorname{Aut}(\operatorname{Mod}(N))\) and/or \(\operatorname{Out}(\operatorname{Mod}(N))\), according to the parenthetical convention. An exact answer must specify the intended one and all allowed \(N\), including small or exceptional topological types, rather than replace the question by a large-genus rigidity assertion.

The summary reports Ivanov-type results for nonorientable mapping class groups and says the full computation for all \(N\) is not settled, but gives no precise range or primary reference. These are supplied status descriptions, not an independently verified all-surface classification.
\subsection{Short English statement}
Compute the automorphism or outer automorphism group of the source's mapping class group \(\operatorname{Mod}(N)\). The surface class and parenthetical choice of target are missing; the summary suggests nonorientable surfaces but does not give precise theorem ranges.
\subsection{Sources}
\begin{itemize}
\item Local UnsolvedMath registry, permanent ID \(11000120\), exact source code \texttt{catalogue entry}, Problem 3.2, PDF page 98. No individual proposer or proposed year is supplied. This page-specific item is distinct from identically numbered problems elsewhere in the volume.
\item Farb, edited volume \emph{Problems on Mapping Class Groups and Related Topics}: \url{https://www.math.uchicago.edu/~farb/papers/mcgbook.pdf}. The volume has 329 explicitly typeset numbered items across chapters, repeated local numbering disambiguated by PDF page.
\item Extraction pdf-font-structure, accessed 2026-07-29. Source evidence NEEDS\_REVIEW presents an open problem and calls for a current primary-literature audit. Supplied OPEN-TRIAGE summary mentions nonorientable Ivanov-type automorphism results without complete citations and says an all-\(N\) computation is not fully settled.
\item Dataset attribution: Ulam AI / UnsolvedMath Contributors, CC BY 4.0. Source rights NEEDS\_REVIEW: publicly accessible author/publisher copy, redistribution terms requiring review.
\end{itemize}
\end{document}
